%% Responsable: Dylber Cabrera
\section{Formalización}
\label{sec:formalizacion}

El problema se formaliza en dos partes, cada una con el mismo esquema: contexto, entradas
y tamaño, salidas, función objetivo y restricciones, y supuestos y alcance. La
demostración de NP-dureza y el análisis de muestreo están en la
sección~\ref{sec:reduccion}; aquí solo se fija el problema que ambas estudian.

\subsection{Notación común}

\begin{definicion}[Espacio de configuraciones]
Hay $L$ categorías y la categoría $i$ ofrece $k_i$ opciones. Una configuración elige una
opción por categoría. El conjunto de configuraciones es $\conf$ y
$|\conf|=\prod_{i=1}^{L}k_i$. Un lote $I$ de $n=|I|$ instancias fija las pruebas: cada
instancia exige un conjunto de restricciones sobre la respuesta del modelo. Una
configuración \emph{cumple} una instancia si la respuesta generada con ella cumple todas
sus restricciones; si falta alguna, no la cumple (puntaje todo-o-nada). Con eso,
$\Acc(c)$ es la fracción de instancias de $I$ que $c$ cumple por completo.
\end{definicion}

Además, $B$ es el presupuesto de generaciones y $m_j$ las consultas gastadas en la
configuración $c_j$. Esta definición de cumplir una instancia no usa el modelo de reglas
de la sección~\ref{sec:reduccion} y es coherente con CONFIG-EXACTA: allí una instancia
se cumple cuando se satisfacen todas las reglas que exige.

\subsection{Parte 1: configuración óptima con oráculo exacto}

\paragraph{Contexto.} Se supone un oráculo que entrega $\Acc(c)$ exacto, sin ruido y sin
límite de consultas. Es una abstracción para estudiar la dureza del problema; el problema
real de esta entrega es la parte 2.

\paragraph{Entrada y tamaño.} La entrada es $(L,k_1,\dots,k_L,I)$. Su tamaño no se reduce
a una sola variable: depende de $L$, de los $k_i$ y de $n$. Como en la
sección~\ref{sec:reduccion}, $L$ y los $k_i$ forman parte de la entrada, de modo que
$|\conf|$ puede crecer de forma exponencial en $L$. El proyecto es el caso $L=5$ y $k_i=8$,
donde $|\conf|=8^5=32\,768$ es una constante: con $L$ fijo el espacio no crece, aunque
siga siendo grande para recorrerlo con un modelo real.

\paragraph{Salida.} Una configuración $c\in\conf$.

\paragraph{Objetivo y restricciones.} Maximizar $\Acc(c)$ sobre $c\in\conf$. No hay
presupuesto de consultas.

\paragraph{Supuestos y alcance.} El oráculo calcula $\Acc(c)$ en tiempo polinomial, como en
CONFIG-EXACTA. Esta parte solo sirve para el análisis de dureza.

\subsection{Parte 2: búsqueda con consultas limitadas o ruidosas}

\paragraph{Contexto.} El oráculo es un modelo de lenguaje y cada consulta es una
generación en GPU. Con temperatura $0{,}0$ la generación es, en principio, determinista
(D2). En GPU pueden aparecer diferencias numéricas pequeñas; esto es una inferencia, no
una medición. El caché hace repetibles las consultas ya hechas. Por eso la variabilidad
relevante no viene del modelo sino de qué instancias se miden, en coherencia con el
análisis de la sección~\ref{sec:reduccion}.

\paragraph{Entrada y tamaño.} El mismo espacio $(L,k_1,\dots,k_L)$, el lote de instancias
$I$, el presupuesto $B$ y la semilla. El tamaño vuelve a depender de $L$, de los $k_i$ y de $n$.

\paragraph{Contabilidad de $B$.} $B$ cuenta las generaciones nuevas por (configuración,
instancia, semilla) (D11). Una pareja ya generada y guardada en el caché no se cobra.

\paragraph{Salida.} Las consultas $m_j$ por configuración $c_j$ y una configuración
ganadora $c^*$.

\paragraph{Objetivo y restricciones.} Maximizar $\Acc(c^*)$ sujeto a $\sum_j m_j\le B$.

\paragraph{Supuestos y alcance.} $c^*$ es la mejor configuración encontrada, no un óptimo
demostrado: el algoritmo no recorre $\conf$ ni garantiza haberlo hecho. Lo que $\Acc(c^*)$
mide es el desempeño sobre el lote de búsqueda, no fuera de él.

\subsection{Planteamiento específico de esta entrega}

\begin{center}
\begin{tabular}{ll}
\toprule
Elemento & Valor \\
\midrule
Espacio explorado & 5 ranuras de 8 opciones, $|\conf|=32\,768$ \\
Temperatura & fija en $0{,}0$ \\
Lote de búsqueda & 100 instancias de \texttt{datos\_visibles.json} \\
Métrica & $\Acc$, todo-o-nada \\
Modelo y semilla & \texttt{qwen17b}, semilla 1 \\
Presupuesto & $B=14\,400$ generaciones (D17) \\
Validación & 87 instancias, una sola vez sobre la configuración elegida \\
\bottomrule
\end{tabular}
\end{center}

La validación se usa solo después de elegir la configuración: si interviniera en la
elección, dejaría de ser una prueba independiente.

\paragraph{Fuera de esta entrega.} Quedan fuera el backtracking, el reparto uniforme
frente a UCB1, los barridos y las gráficas de $T(n)$, y el exacto sobre un subcatálogo
(D17). Successive Halving se descartó y no se implementa (D4).
